\subsection{BLOCK PRODUCTS}

The definitions of no. 4 can be generalized as follows. Let E be a (right or left) A-module, the direct sum of a family \((\mathbf{E}_i)_{i\in I}\) of submodules. For all \(x\in \mathbf{E}\) , let \(x = \sum_{i\in I}x_i\) with \(x_{i}\in \mathbf{E}_{i}\) for all \(i\in I\) ; we shall say that the column matrix \(M(x) = (x_{i})_{i\in I}\) with elements in E is the matrix of \(x\) with respect to the decomposition \((\mathbf{E}_i)_{i\in I}\) of E as a direct sum.

Let F be another A-module (E and F being both right A-modules or both left A-modules) and suppose that F is the direct sum of a family \((\mathbf{F}_{k})_{k\in\mathbb{K}}\) of submodules. For all \(u\in\mathrm{Hom}(\mathbf{E},\mathbf{F})\) and all \(x_{i}\in\mathbf{E}_{i}\) , let \(u(x_{i})=\sum_{k}u_{ki}(x_{i})\) with \(u_{ki}(x_{i})\in\mathbf{F}_{k}\) for all \(k\in K\) ; then \(u_{ki}\in\mathrm{Hom}(\mathbf{E}_{i},\mathbf{F}_{k})\) ; we shall say that the matrix \(M(u)=(u_{ki})_{(k,i)\in K\times I}\) of type (K, I) with elements in the set H the sum of the \(\mathrm{Hom}(\mathbf{E}_{i},\mathbf{F}_{k})\) is the matrix of u with respect to the decompositions \((\mathbf{E}_{i})\) and \((\mathbf{F}_{k})\) of E and F as direct sums.

With these definitions, it is obvious that if \(u, v\) are two A-linear mappings of \(E\) into \(F\) then, for matrices with respect to the same decompositions as direct sums

\[
M (u + v) = M (u) + M (v), \quad M (\gamma u) = \gamma M (u)\tag{21}
\]

for every element \(\gamma\) of the centre of A (no. 2, Remark).

Moreover, the definition of the \(u_{ki}\) shows that, if K is finite, we can write

\[
M (u (x)) = M (u). M (x)\tag{22}
\]

where \(M(u(x))\) is the matrix of \(u(x)\) with respect to the decomposition \((\mathbf{F}_k)\) , the product on the right hand side of (22) being calculated for the mappings \((t, z) \mapsto t(z)\) of \(\operatorname{Hom}(\mathbf{E}_i, \mathbf{F}_k) \times \mathbf{E}_i\) into \(\mathbf{F}_k\) (no. 2, Remark).

Let G be a third A-module, the direct sum of a family \((\mathbf{G}_{l})_{l\in L}\) of submodules, so that there corresponds to every A-linear mapping \(v:F\to G\) a matrix \(M(v)=(v_{lk})\) with respect to the decompositions \((\mathbf{F}_{k})\) and \((\mathbf{G}_{l})\) . If I, K and L are finite, then

\[
M (v \circ u) = M (v). M (u)\tag{23}
\]

where the left hand side is the matrix \((w_{li})\) of \(w = v \circ u\) with respect to the decomposition \((\mathbf{E}_i)\) and \((\mathbf{G}_l)\) and the product on the left hand side is calculated for the mappings \((t, s) \mapsto t \circ s\) of \(\operatorname{Hom}(\mathbf{F}_k, \mathbf{G}_l) \times \operatorname{Hom}(\mathbf{E}_i, \mathbf{F}_k)\) into \(\operatorname{Hom}(\mathbf{E}_i, \mathbf{G}_l)\) (no. 2, Remark). This is just formula (32) of § 1, no. 8, expressed in terms of matrices.

Finally, if I and K are assumed to be finite, \(E^{*}\) (resp. \(F^{*}\) ) is canonically identified with the direct sum of the modules \(E_{i}^{*}\) (resp. \(F_{k}^{*}\) ) ( \(\S\) 2, no. 6, Proposition 10). Then it is immediately verified that the matrix of \(^{t}u\) with respect to the decompositions \((\mathbf{F}_{k}^{*})\) and \((\mathbf{E}_{i}^{*})\) is just \((^{t}u_{kt})_{(k,t)\in\mathbb{K}\times\mathbb{I}}\) .

Suppose now that I and K are finite and further that each of the \(E_{i}\) (resp. \(F_{k}\) ) admits a finite basis. It amounts to the same to say that E (resp. F) admits a basis \((e_{r})_{r\in\mathbb{R}}\) (resp. \((f_{s})_{s\in\mathbb{S}}\) ) and that R (resp. S) admits a partition \((\mathbf{R}_{i})_{i\in\mathbb{I}}\) (resp. \((\mathbf{S}_{k})_{k\in\mathbb{K}}\) ) such that for all \(i\in I\) (resp. \(k\in K\) ), \((e_{r})_{r\in\mathbb{R}_{i}}\) is a basis of \(E_{i}\) (resp. \((f_{s})_{s\in\mathbb{S}_{k}}\) is a basis of \(F_{k}\) ). Then, if \(X=M(u)\) is the matrix of u with respect to the bases \((e_{r})_{r\in\mathbb{R}}\) and \((f_{s})_{s\in\mathbb{S}}\) , the matrix \(X_{ki}=M(u_{ki})\) with respect to the bases \((e_{r})_{r\in\mathbb{R}_{i}}\) and \((f_{s})_{s\in\mathbb{S}_{k}}\) is just the submatrix of X obtained by suppressing the rows of index \(s\notin S_{k}\) and the columns of index \(r\notin R_{i}\) . Thus we define a one-to-one correspondence

\[
X \mapsto (X _ {k i}) _ {(k, i) \in \mathbf {K} \times \mathbf {I}}\tag{24}
\]

between the set of matrices of type (S, R) with elements in A and the set of matrices of matrices \((\mathbf{X}_{ki})_{(k,i)\in\mathbf{K}\times\mathbf{I}}\) of type \(K\times I\) , where each \(X_{ki}\) is a matrix over A of type \((\mathbf{S}_{k},\mathbf{R}_{i})\) . Suppose that further G admits a finite basis \((g_{t})_{t\in\mathbf{T}}\) and that \(\mathbf{T}=(\mathbf{T}_{l})_{l\in\mathbf{L}}\) is a partition of T such that, for each \(l\in L\) , \((g_{t})_{t\in\mathbf{T}_{l}}\) is a basis of \(G_{l}\) ; let \(Y=M(v)\) be the matrix of v with respect to the bases \((f_{s})_{s\in\mathbf{S}}\) and \((g_{t})_{t\in\mathbf{T}}\) , \(Y_{lk}=M(v_{lk})\) that of \(v_{lk}\) with respect to the bases \((f_{s})_{s\in\mathbf{S}_{k}}\) and \((g_{t})_{t\in\mathbf{T}_{l}}\) , Z=M(w) the matrix of w=v∘u with respect to the bases \((e_{r})_{r\in\mathbf{R}}\) and \((g_{t})_{t\in\mathbf{T}}\) and \(Z_{li}=M(w_{li})\) that of \(w_{li}\) with respect to the bases \((e_{r})_{r\in\mathbf{R}_{i}}\) and \((g_{t})_{t\in\mathbf{T}_{l}}\) ; then it follows from (23) that the submatrices \(Z_{li}\) of Z=YX are given by

\[
Z _ {l i} = \sum_ {k} Y _ {l k} X _ {k i}\tag{25}
\]

in other words, the one-to-one correspondence (24) transforms products into products when all the products in question are defined (products of matrices of matrices being defined in the sense of no. 2, Remark); when the submatrices \(Z_{li}\) of the product YX are calculated thus, this product is said to be carried out ``in blocks''.

This name arises from the fact that, when \(I = [1, p]\) and \(K = [1, q]\) , the table representing the matrix X is envisaged as divided into ``blocks'' forming an ``array of matrices''

\[
\left( \begin{array}{c c c c} X _ {1 1} & X _ {1 2} & \ldots & X _ {1 p} \\ X _ {2 1} & X _ {2 2} & \ldots & X _ {2 p} \\ \cdot & \cdot & \cdot & \cdot \\ X _ {q 1} & X _ {q 2} & \ldots & X _ {q p} \end{array} \right)
\]

which is considered as a symbol denoting X when p and q are specific integers sufficiently small for this to be practicable.

